\section{Conclusion}

This study demonstrates that CV models can effectively automate ant counting in behavioral research, addressing the significant limitations of manual counting approaches that have constrained large-scale myrmecological studies. Our findings establish clear guidelines for implementing automated counting systems, showing that CV models with appropriate model calibration strategies can achieve high accuracy across diverse scenarios, from sparse colonies to dense populations exceeding 1,000 individuals per image. While the current approach shows excellent performance within similar environmental conditions to training data, the limited generalizability across different ant species and complex backgrounds highlights important areas for future development, particularly the need for more diverse training datasets and robust data augmentation techniques. The successful detection of pathogen-induced behavioral changes and macronutrient preferences demonstrates the system's capacity to reveal ecological patterns that would be prohibitively difficult to detect through traditional manual methods. By making our complete methodology, datasets, and trained models publicly available through the Ant Detective platform, we provide the myrmecological community with practical tools to conduct larger-scale, more comprehensive behavioral studies that can advance our understanding of social insect ecology and inform broader applications in pest management and ecological monitoring.
